% =====================================================================
% Übungskapitel zu Kapitel 2 (Analytische Geometrie) – steht direkt hinter AnalytischeGeometrie.tex
% Nummerierung: Ü 2.n ; Lösungen im Lösungsteil am Ende des Buches (Loesungen_*.tex).
% =====================================================================
\Uebungskapitel{Übungen zur Analytischen Geometrie}

\Uebungsabschnitt{Vektoren und Vektorraum}

\begin{uebung}*{Vektorketten im Spat}
Die Vektoren $\vv{a}=\vv{AB}$, $\vv{b}=\vv{AD}$ und $\vv{c}=\vv{AE}$ spannen einen Spat $ABCDEFGH$ auf. Die Grundfläche ist das Parallelogramm $ABCD$, die Deckfläche das Parallelogramm $EFGH$, wobei $E$ über $A$, $F$ über $B$, $G$ über $C$ und $H$ über $D$ liegt.
\begin{enumerate}[label=\alph*)]
  \item Stellen Sie die Vektoren $\vv{AG}$, $\vv{BH}$ und $\vv{CE}$ mithilfe der Vektoren $\vv{a}$, $\vv{b}$ und $\vv{c}$ dar.
  \item Der Punkt $M$ ist der Mittelpunkt der Strecke $\overline{FH}$. Stellen Sie den Vektor $\vv{AM}$ mithilfe der Vektoren $\vv{a}$, $\vv{b}$ und $\vv{c}$ dar.
  \item Zeigen Sie mithilfe von Vektorketten, dass sich die Raumdiagonalen $\overline{AG}$ und $\overline{BH}$ gegenseitig halbieren.
\end{enumerate}
\end{uebung}

\begin{uebung}*{Orts- und Verbindungsvektoren}
Gegeben sind die Punkte $A(2|0|1)$, $B(4|2|0)$ und $C(2|4|1)$.
\begin{enumerate}[label=\alph*)]
  \item Berechnen Sie die Koordinaten der Vektoren $\vv{AB}$, $\vv{BC}$ und $\vv{AC}$ sowie deren Längen.
  \item Begründen Sie, dass das Dreieck $ABC$ gleichschenklig ist, und geben Sie die Basis an.
  \item Bestimmen Sie die Koordinaten des Schwerpunkts $S$ des Dreiecks $ABC$.
  \item Bestimmen Sie die Koordinaten des Punktes $D$ so, dass das Viereck $ABCD$ ein Parallelogramm ist. Begründen Sie, dass dieses Parallelogramm eine Raute ist.
\end{enumerate}
\end{uebung}

\Uebungsabschnitt{Skalar-, Vektor- und Spatprodukt}

\begin{uebung}{Skalarprodukt und Winkel}
\begin{enumerate}[label=\alph*)]
  \item Berechnen Sie für das Dreieck $ABC$ aus Ü~2.2 die Größe des Innenwinkels $\beta$ bei $B$.
  \item Untersuchen Sie ohne Taschenrechner, ob das Viereck $ABCD$ aus Ü~2.2 ein Quadrat ist.
  \item Bestimmen Sie ohne Taschenrechner $t\in\mathds{R}$ so, dass die Vektoren $\begin{pmatrix} 1\\ t\\ 2\end{pmatrix}$ und $\begin{pmatrix} 3\\ -1\\ t\end{pmatrix}$ orthogonal zueinander sind.
\end{enumerate}
\end{uebung}

\begin{uebung}*{Vektorprodukt und Spatprodukt}
Gegeben sind die Punkte $O(0|0|0)$, $A(1|0|0)$, $B(0|2|0)$ und $C(0|0|3)$.
\begin{enumerate}[label=\alph*)]
  \item Berechnen Sie $\vv{AB}\times\vv{AC}$ und zeigen Sie, dass der Ergebnisvektor sowohl zu $\vv{AB}$ als auch zu $\vv{AC}$ orthogonal ist.
  \item Berechnen Sie den Flächeninhalt des Dreiecks $ABC$.
  \item Berechnen Sie mithilfe des Spatprodukts das Volumen der dreiseitigen Pyramide $OABC$.
\end{enumerate}
\end{uebung}

\Uebungsabschnitt{Lineare Abhängigkeit und Unabhängigkeit}

\begin{uebung}*{Lineare Abhängigkeit}
\begin{enumerate}[label=\alph*)]
  \item Zeigen Sie, dass die Vektoren $\begin{pmatrix} 1\\ 2\\ 3\end{pmatrix}$, $\begin{pmatrix} 2\\ 3\\ 1\end{pmatrix}$ und $\begin{pmatrix} 3\\ 5\\ 4\end{pmatrix}$ linear abhängig sind.
  \item Bestimmen Sie $a\in\mathds{R}$ so, dass die Vektoren $\begin{pmatrix} 1\\ 0\\ a\end{pmatrix}$, $\begin{pmatrix} 0\\ 1\\ 2\end{pmatrix}$ und $\begin{pmatrix} 1\\ 1\\ 1\end{pmatrix}$ linear abhängig sind.
  \item Erläutern Sie die geometrische Bedeutung der linearen Abhängigkeit dreier Vektoren im Raum.
\end{enumerate}
\end{uebung}

\Uebungsabschnitt{Geraden im Raum}

\begin{uebung}{Spurpunkte, Lagebeziehung und Schnittwinkel von Geraden}
Gegeben sind die Geraden
\[g: \vv{X} = \begin{pmatrix} 1\\ 2\\ 3\end{pmatrix} + \lambda\cdot\begin{pmatrix} 1\\ -1\\ 2\end{pmatrix},\quad
h_1: \vv{X} = \mu\cdot\begin{pmatrix} -2\\ 2\\ -4\end{pmatrix},\]
\[h_2: \vv{X} = \begin{pmatrix} 2\\ 3\\ 7\end{pmatrix} + \mu\cdot\begin{pmatrix} 0\\ 1\\ 1\end{pmatrix},\quad
h_3: \vv{X} = \begin{pmatrix} 0\\ 1\\ 0\end{pmatrix} + \mu\cdot\begin{pmatrix} 1\\ 1\\ 0\end{pmatrix}\]
mit $\lambda,\mu\in\mathds{R}$.
\begin{enumerate}[label=\alph*)]
  \item Bestimmen Sie ohne Taschenrechner die Spurpunkte der Geraden $g$ mit den drei Koordinatenebenen.
  \item Untersuchen Sie ohne Taschenrechner die gegenseitige Lage der Geraden $g$ und $h_1$, $g$ und $h_2$ sowie $g$ und $h_3$. Bestimmen Sie gegebenenfalls die Koordinaten des Schnittpunkts.
  \item Berechnen Sie die Größe des Schnittwinkels der Geraden $g$ und $h_2$.
\end{enumerate}
\end{uebung}

\Uebungsabschnitt{Ebenen im Raum}

\begin{uebung}*{Darstellungsformen von Ebenen}
Die Punkte $A(2|0|0)$, $B(0|4|0)$ und $C(0|0|4)$ legen die Ebene $E$ fest.
\begin{enumerate}[label=\alph*)]
  \item Geben Sie eine Gleichung der Ebene $E$ in Parameterform an.
  \item Bestimmen Sie einen Normalenvektor von $E$ und zeigen Sie, dass $E: 2x_1+x_2+x_3-4=0$ eine Gleichung von $E$ in Koordinatenform ist.
  \item Prüfen Sie, ob der Punkt $P(1|1|1)$ in der Ebene $E$ liegt.
  \item Gegeben ist die Ebene $F: x_1-2x_2+2x_3-6=0$. Bestimmen Sie eine Gleichung von $F$ in Parameterform.
\end{enumerate}
\end{uebung}

\begin{uebung}{Lagebeziehungen und Winkel}
Gegeben sind die Ebenen $E: 2x_1+x_2+x_3-4=0$ und $F: x_1-2x_2+2x_3-6=0$ sowie die Geraden
\[g: \vv{X} = \begin{pmatrix} 1\\ 1\\ 4\end{pmatrix} + \lambda\cdot\begin{pmatrix} 1\\ -1\\ -1\end{pmatrix}\quad\text{und}\quad
k: \vv{X} = \begin{pmatrix} 0\\ 0\\ 1\end{pmatrix} + \lambda\cdot\begin{pmatrix} 1\\ 1\\ 1\end{pmatrix}\quad(\lambda\in\mathds{R}).\]
\begin{enumerate}[label=\alph*)]
  \item Zeigen Sie ohne Taschenrechner, dass die Gerade $g$ echt parallel zur Ebene $E$ verläuft.
  \item Bestimmen Sie ohne Taschenrechner die Koordinaten des Schnittpunkts $S$ der Geraden $k$ mit der Ebene $E$ und berechnen Sie die Größe des Schnittwinkels zwischen $k$ und $E$.
  \item Begründen Sie, dass sich die Ebenen $E$ und $F$ schneiden, und bestimmen Sie eine Gleichung der Schnittgeraden.
  \item Berechnen Sie die Größe des Schnittwinkels der Ebenen $E$ und $F$.
\end{enumerate}
\end{uebung}

\Uebungsabschnitt{Abstände im Raum}

\begin{uebung}*{Abstände}
\begin{enumerate}[label=\alph*)]
  \item Berechnen Sie den Abstand des Punktes $P(3|4|5)$ von der Ebene $E: x_1+2x_2+2x_3-6=0$.
  \item Gegeben ist die Gerade \[g: \vv{X} = \begin{pmatrix} 1\\ 2\\ 3\end{pmatrix} + \lambda\cdot\begin{pmatrix} 1\\ -1\\ 2\end{pmatrix}\quad(\lambda\in\mathds{R}).\] Bestimmen Sie den Abstand des Punktes $Q(3|0|4)$ von der Geraden $g$ sowie die Koordinaten des zugehörigen Lotfußpunkts.
  \item Gegeben ist zusätzlich die Gerade \[h: \vv{X} = \begin{pmatrix} 0\\ 1\\ 0\end{pmatrix} + \mu\cdot\begin{pmatrix} 1\\ 1\\ 0\end{pmatrix}\quad(\mu\in\mathds{R}).\] Berechnen Sie den Abstand der windschiefen Geraden $g$ und $h$.
\end{enumerate}
\end{uebung}

\Uebungsabschnitt{Anwendungen}

\begin{uebung}*{Kugel und Ebene}
Gegeben sind die Kugel $K: (x_1-2)^2+(x_2-1)^2+(x_3+1)^2 = 25$ und die Ebene $E: x_1+2x_2+2x_3+7=0$.
\begin{enumerate}[label=\alph*)]
  \item Geben Sie den Mittelpunkt $M$ und den Radius $R$ der Kugel an. Prüfen Sie, ob der Punkt $P(5|5|-1)$ auf der Kugel liegt und ob der Punkt $Q(4|3|0)$ innerhalb oder außerhalb der Kugel liegt.
  \item Zeigen Sie, dass die Ebene $E$ die Kugel $K$ in einem Kreis schneidet.
  \item Bestimmen Sie den Mittelpunkt und den Radius des Schnittkreises.
\end{enumerate}
\end{uebung}

\begin{uebung}*{Spiegelung eines Punktes an einer Ebene}
Gegeben sind die Ebene $E: x_1+2x_2+2x_3-6=0$ und der Punkt $P(3|3|3)$.
\begin{enumerate}[label=\alph*)]
  \item Geben Sie eine Gleichung der Lotgeraden $\ell$ zu $E$ durch $P$ an.
  \item Bestimmen Sie die Koordinaten des Lotfußpunkts $F$ und den Abstand des Punktes $P$ von der Ebene $E$.
  \item Bestimmen Sie die Koordinaten des Bildpunkts $P'$ bei der Spiegelung von $P$ an der Ebene $E$.
\end{enumerate}
\end{uebung}

\begin{uebung}{Pyramide (Abiturniveau)}
Gegeben sind die Punkte $A(0|0|0)$, $B(4|0|0)$, $C(4|4|0)$, $D(0|4|0)$ und $S(2|2|6)$. Sie sind die Eckpunkte einer Pyramide $ABCDS$ mit der Spitze $S$.
\begin{enumerate}[label=\alph*)]
  \item Zeigen Sie, dass die Grundfläche $ABCD$ ein Quadrat ist, und berechnen Sie das Volumen der Pyramide.
  \item Bestimmen Sie eine Gleichung der Ebene $L$, in der die Seitenfläche $BCS$ liegt, in Koordinatenform.\\
  {\small (zur Kontrolle: $L: 3x_1+x_3-12=0$)}
  \item Berechnen Sie die Größe des Winkels, den die Seitenfläche $BCS$ mit der Grundfläche einschließt.
  \item Berechnen Sie den Abstand des Punktes $A$ von der Ebene $L$.
  \item Der Punkt $N$ ist der Mittelpunkt der Kante $\overline{BC}$. Berechnen Sie die Länge der Strecke $\overline{NS}$ und damit den Flächeninhalt des Dreiecks $BCS$. Bestätigen Sie Ihr Ergebnis mithilfe des Vektorprodukts.
\end{enumerate}
\end{uebung}

\Uebungsabschnitt{Abiturniveau \textendash\ Teil A (ohne Hilfsmittel)}

\begin{uebung}*{Ebene und Abstand}
Gegeben ist die Ebene $E: 2x_1-x_2+2x_3-6=0$.
\begin{enumerate}[label=\alph*)]
  \item Bestimmen Sie die Koordinaten der Schnittpunkte von $E$ mit den Koordinatenachsen.
  \item Berechnen Sie den Abstand des Koordinatenursprungs von $E$.
  \item Geben Sie die Gleichungen aller Ebenen an, die parallel zu $E$ sind und von $E$ den Abstand $2$ haben.
\end{enumerate}
\end{uebung}

\begin{uebung}*{Gerade und Kugel}
Gegeben sind die Kugel $K$ mit dem Mittelpunkt $M(1|2|-1)$ und dem Radius $3$ sowie die Gerade
\[g: \vv{X} = \begin{pmatrix} 1\\ 2\\ -1\end{pmatrix}+\lambda\cdot\begin{pmatrix} 2\\ 1\\ 2\end{pmatrix}\quad(\lambda\in\mathds{R}).\]
\begin{enumerate}[label=\alph*)]
  \item Begründen Sie, dass $g$ durch den Mittelpunkt von $K$ verläuft, und bestimmen Sie die Koordinaten der beiden Schnittpunkte von $g$ und $K$.
  \item Im Schnittpunkt $P$ mit positivem Parameterwert wird die Tangentialebene $T$ an $K$ gelegt. Bestimmen Sie eine Gleichung von $T$ in Koordinatenform.
\end{enumerate}
\end{uebung}

\begin{uebung}*{Gerade und Ebene mit Parameter}
Gegeben sind die Ebene $E: x_1+2x_2-2x_3-4=0$ und für $a\in\mathds{R}$ die Geraden
\[g_a: \vv{X} = \begin{pmatrix} 0\\ 0\\ a\end{pmatrix}+\lambda\cdot\begin{pmatrix} 2\\ -1\\ 0\end{pmatrix}\quad(\lambda\in\mathds{R}).\]
\begin{enumerate}[label=\alph*)]
  \item Zeigen Sie, dass jede Gerade $g_a$ parallel zu $E$ verläuft.
  \item Bestimmen Sie $a$ so, dass $g_a$ in $E$ liegt.
\end{enumerate}
\end{uebung}

\begin{uebung}*{Rechtwinkliges Dreieck}
Gegeben sind die Punkte $A(1|0|2)$, $B(3|2|1)$ und $C(3|-1|4)$.
\begin{enumerate}[label=\alph*)]
  \item Zeigen Sie, dass das Dreieck $ABC$ bei $A$ rechtwinklig und gleichschenklig ist.
  \item Berechnen Sie den Flächeninhalt des Dreiecks.
  \item Bestimmen Sie die Koordinaten des Punktes $D$ so, dass $ABDC$ ein Quadrat ist.
\end{enumerate}
\end{uebung}

\Uebungsabschnitt{Abiturniveau \textendash\ Teil B (mit Hilfsmitteln)}

\begin{uebung}{Dachfläche mit Antenne}
In einem kartesischen Koordinatensystem (Längeneinheit $1\,\text{m}$) beschreibt die $x_1x_2$-Ebene den horizontalen Boden. Eine ebene Dachfläche enthält die Punkte $P(0|0|3)$, $Q(10|0|3)$ und $R(0|8|9)$.
\begin{enumerate}[label=\alph*)]
  \item Bestimmen Sie eine Gleichung der Ebene $D$, in der die Dachfläche liegt, in Koordinatenform.\\
  {\small (zur Kontrolle: $D: -3x_2+4x_3-12=0$)}
  \item Berechnen Sie den Neigungswinkel der Dachfläche gegenüber dem Boden.
  \item Im Punkt $F(5|4|6)$ der Dachfläche steht eine senkrechte (d.\,h. zur $x_3$-Achse parallele) Antenne der Länge $2\,\text{m}$ mit der Spitze $S$. Zeigen Sie, dass $F$ in $D$ liegt, und berechnen Sie den Abstand der Antennenspitze von der Dachfläche.
  \item Sonnenlicht fällt in Richtung des Vektors $\vv{v} = \begin{pmatrix} 1\\ 0\\ -2\end{pmatrix}$ ein. Bestimmen Sie den Schattenpunkt der Antennenspitze auf der Dachfläche und die Länge des Schattens der Antenne.
\end{enumerate}
\end{uebung}

\begin{uebung}{Flugbahnen}
Zwei Flugzeuge bewegen sich geradlinig mit konstanter Geschwindigkeit. Ihre Positionen werden für $t\geq 0$ beschrieben durch
\[F_1: \vv{X}(t) = \begin{pmatrix} 0\\ 0\\ 8\end{pmatrix}+t\cdot\begin{pmatrix} 4\\ 3\\ 0\end{pmatrix},\qquad F_2: \vv{X}(t) = \begin{pmatrix} 24\\ 0\\ 7\end{pmatrix}+t\cdot\begin{pmatrix} -4\\ 3\\ 0\end{pmatrix}\]
($t$ in Minuten ab Beobachtungsbeginn, Koordinaten in km, die $x_3$-Achse zeigt senkrecht nach oben).
\begin{enumerate}[label=\alph*)]
  \item Berechnen Sie die Geschwindigkeiten der Flugzeuge in $\frac{\text{km}}{\text{h}}$.
  \item Zeigen Sie, dass die Flugbahnen windschief sind.
  \item Berechnen Sie den Abstand der beiden Flugbahnen.
  \item Bestimmen Sie den Abstand $d(t)$ der Flugzeuge zum Zeitpunkt $t$ sowie den Zeitpunkt, zu dem die Flugzeuge einander am nächsten kommen. Vergleichen Sie mit c).
  \item Berechnen Sie die Größe des Winkels, unter dem sich die Flugrichtungen (in der Draufsicht) kreuzen.
\end{enumerate}
\end{uebung}

\begin{uebung}{Schnitt durch einen Würfel}
Der Würfel $ABCDEFGH$ hat die Eckpunkte $A(0|0|0)$, $B(6|0|0)$, $D(0|6|0)$, $E(0|0|6)$ und $G(6|6|6)$.
\begin{enumerate}[label=\alph*)]
  \item Die Ebene $L$ enthält die Punkte $B$, $D$ und $E$. Bestimmen Sie eine Gleichung von $L$ in Koordinatenform.
  \item Berechnen Sie die Abstände der Punkte $A$ und $G$ von $L$.
  \item Zeigen Sie, dass die Raumdiagonale $\overline{AG}$ senkrecht auf $L$ steht, und bestimmen Sie, in welchem Verhältnis $L$ diese Raumdiagonale teilt.
  \item $L$ schneidet von dem Würfel eine dreiseitige Pyramide mit der Spitze $A$ ab. Berechnen Sie, welchen Anteil ihr Volumen am Würfelvolumen hat.
\end{enumerate}
\end{uebung}

\begin{uebung}{Kugel auf einer schiefen Ebene}
Eine Kugel $K$ mit dem Mittelpunkt $M(4|1|4)$ und dem Radius $3$ liegt an der Ebene $E: 2x_1-x_2+2x_3-6=0$ an.
\begin{enumerate}[label=\alph*)]
  \item Zeigen Sie, dass $E$ die Kugel berührt, und bestimmen Sie die Koordinaten des Berührpunkts $B$.
  \item Bestimmen Sie eine Gleichung der zweiten zu $E$ parallelen Ebene, die $K$ berührt.
  \item Geben Sie die Koordinaten des höchsten Punkts der Kugel an (die $x_3$-Achse zeigt nach oben).
\end{enumerate}
\end{uebung}

\begin{uebung}{Zelt}
Ein Zelt hat als Boden das Rechteck mit den Eckpunkten $A(0|0|0)$, $B(4|0|0)$, $C(4|6|0)$ und $D(0|6|0)$. Der First verläuft von $F(2|0|2)$ nach $G(2|6|2)$ (Angaben in m). Die beiden Seitenwände sind die Rechtecke $ADGF$ und $BCGF$.
\begin{enumerate}[label=\alph*)]
  \item Bestimmen Sie Gleichungen der Ebenen $W_1$ (durch $A$, $D$, $F$) und $W_2$ (durch $B$, $C$, $F$) in Koordinatenform.
  \item Berechnen Sie den Winkel zwischen den beiden Seitenwänden.
  \item Berechnen Sie das Volumen des Zeltes.
  \item Im Punkt $P(1|3|0)$ wird eine senkrechte Stange aufgestellt. Bestimmen Sie die Höhe, in der sie die Zeltwand berührt.
  \item Im Punkt $L(2|3|1)$ hängt eine Lampe. Berechnen Sie ihren Abstand von der Seitenwand $W_1$.
\end{enumerate}
\end{uebung}

\begin{uebung}{Reflektierter Laserstrahl}
Ein Laserstrahl geht vom Punkt $Q(0|0|6)$ aus und verläuft in Richtung des Vektors $\vv{u} = \begin{pmatrix} 1\\ 1\\ -2\end{pmatrix}$. Er trifft auf den Boden ($x_1x_2$-Ebene), wird dort nach dem Reflexionsgesetz gespiegelt und trifft dann auf eine Wand in der Ebene $x_1 = 6$.
\begin{enumerate}[label=\alph*)]
  \item Bestimmen Sie den Auftreffpunkt $R$ auf dem Boden.
  \item Begründen Sie, dass der reflektierte Strahl die Richtung $\vv{u'} = \begin{pmatrix} 1\\ 1\\ 2\end{pmatrix}$ hat, und bestimmen Sie den Auftreffpunkt $W$ auf der Wand.
  \item Berechnen Sie den Winkel, unter dem der Strahl auf den Boden trifft.
\end{enumerate}
\end{uebung}

\UebungskapitelEnde
